% !TEX root = ../main.tex
% !TEX program = lualatex
\documentclass[../main.tex]{subfiles}
\IfSubfilesClassLoaded{\externaldocument{\subfix{../build/main}}}{}
% ====================
% File: 18A_Ship_Trials.tex
% ====================
\begin{document}
\IfSubfilesClassLoaded{\chapter{Ship Trials}}{}
% === Ship Trials (EDO 3.5.7) ===
\section{Ship Trials}

\subsection{Summary}
\begin{description}
  \item[\textbf{Purpose.}] Ship trials verify readiness, safety, and contractual compliance before delivery and Fleet introduction~\autocite{edo-3-5-7-ship-trials-2025}.
  \item[\textbf{Major trial types.}] Surface ship trials include \ac{bdt}, \ac{bst}, \ac{at}, \ac{iat}, \ac{pdtt}, \ac{fct}, \ac{st}, and post-shakedown availability events~\autocite{edo-3-5-7-ship-trials-2025}.
  \item[\textbf{Key players.}] \ac{insurv}, \ac{opnav}, \ac{navsea}, \ac{peo}, \ac{tycom}, \ac{supship}, and ship's force each hold defined roles in trials and acceptance~\autocite{edo-3-5-7-ship-trials-2025}.
  \item[\textbf{Submarine emphasis.}] Submarine trials include fast cruise, ALPHA, BRAVO, \ac{insurv} combined trials, and \ac{gmi}, with SUBSAFE certification and \ac{uro} authorization~\autocite{edo-3-5-7-ship-trials-2025}.
  \item[\textbf{Certification.}] Post-sea trial certification requires completion messages, deficiency tracking, and readiness determinations for Fleet operations~\autocite{edo-3-5-7-ship-trials-2025}.
\end{description}

\subsection{Ship Trial Requirements}
\begin{description}
  \item[\textbf{Title 10.}] 10~U.S.C.~8674 directs \ac{secnav} to designate boards of naval officers to examine naval vessels, including unfinished vessels, and treats vessels undergoing necessary trials before fleet acceptance as a distinct case~\autocite{USC-10-8674,edo-3-5-7-ship-trials-2025}.
  \item[\textbf{OPNAV policy.}] \ac{opnavinst}~4700.8M governs trials, acceptance, commissioning, fitting out, shakedown, and post-shakedown availability for ships under construction or conversion~\autocite{OPNAVINST-4700-8M,edo-3-5-7-ship-trials-2025}.
  \item[\textbf{Mission capable prior to delivery.}] Ships should be fully mission capable before delivery, except for crew certification or range requirements not possible until after delivery~\autocite{edo-3-5-7-ship-trials-2025}.
\end{description}

\subsection{Key Terms}
\begin{description}
  \item[\textbf{Accepting authority.}] Officer designated by the \ac{cno} to accept a vessel for the Navy, normally the cognizant \ac{peo}~\autocite{edo-3-5-7-ship-trials-2025}.
  \item[\textbf{Presenting authority.}] Officer designated to present the ship to the trial board and certify readiness for trials~\autocite{edo-3-5-7-ship-trials-2025}.
  \item[\textbf{OWLD.}] Obligations for construction and post-delivery work must occur before the \ac{owld}, typically 11 months after completion of fitting out~\autocite{edo-3-5-7-ship-trials-2025}.
  \item[\textbf{Delivery and guarantee period.}] Delivery is the date the Navy accepts the ship; the shipbuilder remains responsible for correcting defects during the guarantee period~\autocite{edo-3-5-7-ship-trials-2025}.
\end{description}

\subsection{Key Players and Responsibilities}
\begin{itemize}
  \item \ac{insurv} provides independent verification of readiness for acceptance and Fleet introduction~\autocite{edo-3-5-7-ship-trials-2025}.
  \item \ac{opnav} makes final determinations of readiness for service; \ac{navsea} provides formal delivery recommendations to \ac{cno} via \ac{asnrdanda}~\autocite{edo-3-5-7-ship-trials-2025}.
  \item \ac{peo} acts as accepting authority for new construction and conversion ships; \ac{supship} prepares, presents, and documents discrepancies~\autocite{edo-3-5-7-ship-trials-2025}.
  \item \ac{tycom} represents the warfighter; Fleet commanders provide trial services; ship's force executes post-delivery trials~\autocite{edo-3-5-7-ship-trials-2025}.
  \item The Naval Supervising Authority assumes custody of the ship's material condition during industrial periods~\autocite{edo-3-5-7-ship-trials-2025}.
\end{itemize}

\subsection{Surface Ship Construction Trials}
\begin{description}
  \item[\textbf{BDT.}] Pier-side tests to determine readiness to safely conduct sea trials and rehearse the \ac{at} schedule of events; presented by shipbuilder to \ac{supship}~\autocite{edo-3-5-7-ship-trials-2025}.
  \item[\textbf{BST.}] Conducted after \ac{bdt} to demonstrate seaworthiness and readiness for \ac{at}; presented by shipbuilder to \ac{supship} with \ac{navsea} and ship's force observers~\autocite{edo-3-5-7-ship-trials-2025}.
  \item[\textbf{AT.}] Conducted after \ac{bst} to assess contract compliance and operational readiness; \ac{insurv} documents discrepancies and recommends acceptance or additional trials~\autocite{edo-3-5-7-ship-trials-2025}.
  \item[\textbf{IAT.}] Consolidates \ac{bst} and \ac{at} before delivery to accelerate schedules for mature production lines~\autocite{edo-3-5-7-ship-trials-2025}.
\end{description}

\subsection{Post Delivery Test and Trials}
\begin{description}
  \item[\textbf{When.}] Conducted prior to the shakedown period~\autocite{edo-3-5-7-ship-trials-2025}.
  \item[\textbf{Purpose.}] Executes warfare-area trials such as acoustic trials, propulsion plant examinations, \ac{cssqt} (surface) or weapons system accuracy trials (submarines), and ASW/weapon accuracy events~\autocite{edo-3-5-7-ship-trials-2025}.
  \item[\textbf{Responsibility.}] Conducted by the ship's commanding officer and ship's force, with special assistance teams as needed~\autocite{edo-3-5-7-ship-trials-2025}.
\end{description}

\subsection{Final Contract Trials}
\begin{description}
  \item[\textbf{When.}] After \ac{at} deficiency corrections and prior to post-shakedown availability; before guarantee period expiration~\autocite{edo-3-5-7-ship-trials-2025}.
  \item[\textbf{Purpose.}] Final check for shipbuilder-related defects and operational readiness; similar in scope to \ac{at} but executed by Navy personnel~\autocite{edo-3-5-7-ship-trials-2025}.
  \item[\textbf{Responsibility.}] Presented by the shipbuilding program manager to \ac{insurv} and observed by \ac{navsea}, \ac{supship}, and the shipbuilder~\autocite{edo-3-5-7-ship-trials-2025}.
\end{description}

\subsection{Shock Trials}
\begin{description}
  \item[\textbf{When.}] Conducted prior to post-shakedown availability for first-of-class ships~\autocite{edo-3-5-7-ship-trials-2025}.
  \item[\textbf{Purpose.}] Demonstrate the ability to fight the ship and satisfy Congressional \ac{lfte} requirements for covered systems when applicable~\autocite{USC-10-4172,DoDI5000-98,edo-3-5-7-ship-trials-2025}.
  \item[\textbf{Responsibility.}] Shipbuilding program manager oversees; \ac{nswc} Carderock designs and instruments trials; ship's CO serves as Officer in Tactical Command~\autocite{edo-3-5-7-ship-trials-2025}.
\end{description}

\subsection{Post-Shakedown Availability}
\begin{description}
  \item[\textbf{When.}] After shakedown and prior to expiration of \ac{scn} funding and the \ac{owld}~\autocite{edo-3-5-7-ship-trials-2025}.
  \item[\textbf{Purpose.}] Correct new construction deficiencies and accomplish authorized improvements~\autocite{edo-3-5-7-ship-trials-2025}.
  \item[\textbf{Responsibility.}] Supervising authority and ship's CO determine test scope; ship's force operates systems and the shipyard remains responsible for completion~\autocite{edo-3-5-7-ship-trials-2025}.
\end{description}

\subsection{Special Trials}
\begin{description}
  \item[\textbf{NAVSEA trials.}] Conducted on representative ships after delivery to determine class characteristics or for ship-specific needs~\autocite{edo-3-5-7-ship-trials-2025}.
  \item[\textbf{Examples.}] Standardization, tactical, acoustical, vibration, fuel economy, heat balance, emergency recovery, and stability/control trials~\autocite{edo-3-5-7-ship-trials-2025}.
\end{description}

\subsection{Submarine Trials Timeline (New Construction)}
\begin{description}
  \item[\textbf{Fast cruise.}] \ac{tycom}-directed evolutions for 96 hours with required crew rest to support certification~\autocite{edo-3-5-7-ship-trials-2025}.
  \item[\textbf{ALPHA trials.}] Propulsion plant full-power trials, test depth, ship control, and initial tightness checks~\autocite{edo-3-5-7-ship-trials-2025}.
  \item[\textbf{BRAVO trials.}] Combat systems, navigation, communications, and acoustic performance to prepare for \ac{insurv} trials~\autocite{edo-3-5-7-ship-trials-2025}.
  \item[\textbf{BRAVO work period.}] Correct discrepancies and conduct post-dive checks and inspections~\autocite{edo-3-5-7-ship-trials-2025}.
  \item[\textbf{INSURV trials and delivery.}] \ac{insurv} conducts combined acceptance and final contract trials; delivery is followed by \ac{gmi}~\autocite{edo-3-5-7-ship-trials-2025}.
\end{description}

\subsection{Post-Industrial Availability Trials}
\begin{description}
  \item[\textbf{Purpose.}] Conducted after repairs to verify readiness for service; scope depends on the work package~\autocite{edo-3-5-7-ship-trials-2025}.
  \item[\textbf{Full power.}] Required after each regular overhaul; other repairs use trial scopes agreed by oversight and the ship's CO~\autocite{edo-3-5-7-ship-trials-2025}.
\end{description}

\subsection{Post-Sea Trials Certification}
\begin{description}
  \item[\textbf{CNO availabilities.}] Naval supervising authority, lead maintenance activity, and ship's force conduct a departure conference and issue an availability completion message detailing unresolved work and configuration changes~\autocite{edo-3-5-7-ship-trials-2025}.
  \item[\textbf{Submarine requirements.}] \ac{navsea} certifies SUBSAFE material condition and recommends \ac{uro}; \ac{tycom} authorizes \ac{uro} to test depth~\autocite{edo-3-5-7-ship-trials-2025}.
\end{description}

\subsection{Submarine New Construction vs. Repair}
\begin{description}
  \item[\textbf{NAVSEA certification.}] New construction certifies SUBSAFE material condition for the entire submarine; repair certifies only the work performed~\autocite{edo-3-5-7-ship-trials-2025}.
  \item[\textbf{Testing.}] New construction uses multiple at-sea trial periods; repair typically uses one trial period to complete all testing~\autocite{edo-3-5-7-ship-trials-2025}.
\end{description}

\subsection{Ship Program OT\&E}
\begin{description}
  \item[\textbf{Dedicated OT phases.}] Per COMOPTEVFORINST~3980.2, ships use EOA, OA, \ac{iotre}, VCD, and \ac{fote} to assess operational suitability and effectiveness~\autocite{edo-3-5-7-ship-trials-2025}.
  \item[\textbf{Program coordination.}] \ac{mdauthority} and \ac{dote} align on test requirements; \ac{dte} and \ac{ote} may be combined with \ac{secnav} concurrence, but \ac{techeval} and \ac{iote} are not combined~\autocite{edo-3-5-7-ship-trials-2025}.
  \item[\textbf{Pre-Milestone B testing.}] \ac{dte} and \ac{ote} before Milestone~B address individual systems, land-based testing, or validated modeling and simulation when necessary~\autocite{edo-3-5-7-ship-trials-2025}.
\end{description}

\subsection{Coursebook cue: ship-trial types and roles}

Coursebook 3.5.7 expects more than the names of trials.  The board answer
should connect timing, lead organization, certification purpose, and platform
distinctions
\autocite{edo-3-5-7-ship-trials-2025}.

\begin{longtblr}[
  caption = {Ship-Trial Board Cue},
  label = {tab:ship-trial-board-cue},
]{
  width = \linewidth,
  colspec = {X[0.75,l] X[1.8,l]},
  rowhead = 1,
  hlines,
  vlines,
}
\textbf{Trial cue} & \textbf{What to say} \\
Builder's trials & Contractor-led event used to prove readiness before government acceptance events. \\
Acceptance trials & Government-witnessed event used to support delivery or acceptance decisions. \\
Final contract trials & Post-delivery event used to verify correction of deficiencies and contract compliance. \\
Special trials & Focused events for unique systems, mission areas, or certification needs. \\
Role split & Builder, supervisor of shipbuilding, INSURV, program office, TYCOM, and technical authorities each have distinct interests. \\
Submarine distinction & Nuclear, SUBSAFE, deep-submergence, and certification constraints can change trial sequencing and evidence requirements. \\
\end{longtblr}

\subsection{Quick Review}
\begin{itemize}
  \item Ship trials are anchored by Title~10 and \ac{opnavinst} policy to support material-readiness, acceptance, and delivery decisions~\autocite{USC-10-8674,OPNAVINST-4700-8M,edo-3-5-7-ship-trials-2025}.
  \item \ac{bdt}, \ac{bst}, \ac{at}, and \ac{iat} are the core surface-ship construction trials; \ac{pdtt} and \ac{fct} follow delivery~\autocite{edo-3-5-7-ship-trials-2025}.
  \item Submarine trials culminate in \ac{insurv} combined trials, \ac{gmi}, and SUBSAFE certification with \ac{uro} authorization~\autocite{edo-3-5-7-ship-trials-2025}.
  \item Post-sea trial certification and availability completion messages confirm readiness and track deficiencies~\autocite{edo-3-5-7-ship-trials-2025}.
\end{itemize}

%====================
% End of file
%====================
\IfSubfilesClassLoaded{
  \RenewDocumentCommand{\entryname}{}{\textbf{\color{Modern} Acronym}}
  \RenewDocumentCommand{\descriptionname}{}{\textbf{\color{Modern} Definition}}
  \printnoidxglossary[
    type=\acronymtype,
    title=Acronyms,
    style=long-booktabs
  ]}{}

\end{document}
